\documentclass[12pt]{article}

\usepackage[utf8]{inputenc}
\usepackage[T1]{fontenc}
\usepackage[margin=1in]{geometry}
\usepackage{parskip}
\usepackage{enumitem}
\usepackage{graphicx}
\usepackage{hyperref}

\title{Engineering for Clarity: Designing SimplrAI,\\
a Grade-Appropriate AI Explanation System for Student Learning}
\author{Soumil Bhandari \\ BASIS Independent McLean}
\date{April 2026}

\begin{document}
\maketitle

\section{Introduction}

Over the past two years, artificial intelligence has moved from a futuristic
concept to a core academic tool. According to the Pew Research Center (2024),
more than 63\% of American teenagers report using AI for schoolwork, and nearly
one in five uses it daily. The tools they reach for---ChatGPT, Gemini,
Perplexity---were built as general-purpose conversational agents, optimized
for breadth, fluency, and general reasoning. None were designed for the specific
cognitive needs of a thirteen-year-old trying to understand the chain rule for
the first time.

In my own coursework, and during the year I spent tutoring middle- and
high-school students, I encountered the same pattern repeatedly. A student would
ask a chatbot to explain mitosis, derivative notation, or the causes of the
French Revolution. The model would produce a fluent, dense, six-paragraph
response written somewhere between a college textbook and an encyclopedia
article. The student would skim it, decide they were still confused, and ask
for it to be simpler. The next response would often be longer than the first.

This is a design problem, not a capability problem. Frontier LLMs already have
the underlying knowledge to teach high-school content; what they lack is a
pedagogical scaffold---structured output, age-aware reading level, intentional
hallucination control, and an interface that respects what cognitive scientists
have known for forty years about working-memory limits (Sweller, 1988;
Paas \& van Merri\"enboer, 2020).

This project asks:

\begin{quote}\itshape
How can an AI-based academic explanation system be engineered to produce fast,
accurate, and grade-appropriate explanations that reduce cognitive load and
meaningfully improve student comprehension?
\end{quote}

To answer it, I designed and built SimplrAI, an explanation system for students
from sixth grade through the first years of college. SimplrAI wraps a frontier
LLM in three deliberate layers:
\begin{enumerate}[leftmargin=1.4em]
  \item A \textbf{structured prompting layer} enforcing a
  definition--example--analogy--application output schema with lightweight
  retrieval grounding.
  \item A \textbf{grade-level adaptation layer} matching vocabulary and sentence
  complexity to the target Lexile band.
  \item A \textbf{minimal user interface} built on progressive-disclosure
  principles.
\end{enumerate}

Section 2 surveys the relevant literature. Section 3 describes methodology.
Sections 4--6 present results, discussion, and limitations.

\section{Literature Review}

\subsection{AI tutoring systems}

Koedinger et al.\ (2022) found that AI tutors can support conceptual
understanding when their explanations are appropriately structured and
scaffolded. OpenAI (2023) and Anthropic (2024) both describe their models as
optimized for general helpfulness rather than pedagogical clarity. Pew Research
(2024) confirmed rapid adoption but warned educational usefulness is limited by
inconsistent accuracy. Dwivedi et al.\ (2023) argue LLMs produce outputs that
are syntactically fluent but epistemically unreliable.

\subsection{Cognitive load theory}

Cognitive load theory (Sweller, 1988) holds that working memory can only process
a limited amount of novel information at once. Long, dense, disorganized
explanations increase extraneous load. Later refinements (Kalyuga, 2009;
Paas \& van Merri\"enboer, 2020) emphasize concise, sequenced, chunked
explanations. Chi and Wylie's ICAP framework (2014) finds explanations prompting
active construction outperform passive ones.

\subsection{Hallucination and factual reliability}

Maynez et al.\ (2020) document hallucination rigorously. OpenAI (2023) shows
prompting strategy influences accuracy. Mitigation approaches: structured
prompting (Li et al., 2023) and retrieval-augmented generation (Lewis et al.,
2020). SimplrAI combines both.

\subsection{Reading-level adaptation and interface design}

The Lexile Framework (MetaMetrics, 2019) shows comprehension falls when text
exceeds a learner's threshold. Xu et al.\ (2023) show LLMs default to
college-level prose unless constrained. Nielsen Norman Group (2023), Reich
(2020), and Quinn et al.\ (2023) find students learn best from interfaces with
low visual noise and chunked content.

\subsection{Synthesis and gaps}

Three gaps: no deployed system targets grade-appropriate explanation; structured
prompts unevaluated for student comprehension; interface effects understudied.
SimplrAI addresses all three.

\section{Methodology}

\subsection{The SimplrAI system}

SimplrAI is a React/Node/Express/Supabase web application deployed at
\href{https://simplrai.tech}{\texttt{simplrai.tech}}. The \textit{prompting
layer} enforces a four-part output (definition, example, analogy, application),
imposes a maximum response length, and forbids above-grade terminology without
inline definition. The \textit{retrieval layer} searches a 4,300-document corpus
from OpenStax, Khan Academy, and CK-12; up to three passages are passed as
grounding context. The \textit{interface layer} renders each section as an
expandable card with a ``explain simpler'' button.

\subsection{Participants}

Students were recruited from two U.S.\ partner schools and an online
study-skills community. Final sample: $N = 73$ (middle school $n = 24$;
high school $n = 31$; college $n = 18$). Ages 11--22 ($M = 16.2$, $SD = 2.7$);
49\% female, 47\% male. All minors provided parental consent.

% NOTE: double check N. i think a few more responses came in
% might need to re-run exclusions

\subsection{Design}

Three-arm between-subjects design: SimplrAI, ChatGPT (GPT-5), textbook-only
control. Stratified randomization balanced grade level. Each participant
completed Pre-Algebra, Biology, and U.S.\ History sessions (counterbalanced).

\begin{figure}[h]
  \centering
  \includegraphics[width=\linewidth]{fig1_design.pdf}
  \caption{Study design and participant flow.}
  \label{fig:design}
\end{figure}

\subsection{Procedure}

10-question pre-test $\to$ 20 minutes with assigned tool $\to$ 10-question
post-test $\to$ 5-item Likert survey (clarity, accuracy, helpfulness, reading
difficulty, trust).

\subsection{Materials and measures}

Quiz items validated by two teachers and pilot-tested with five students.
Reading level measured via Flesch--Kincaid. Factual accuracy scored by two
experts on a four-point rubric; weighted Cohen's $\kappa = 0.84$.

\subsection{Analysis plan}

Descriptive statistics; independent-samples $t$-tests; Cohen's $d$.
All planned comparisons reported (pre-registered).

\end{document}
